                              IMC 2016, Blagoevgrad, Bulgaria


                                               Day 2, July 28, 2016




                                                                                                               ∞
                                                                                                               X     xn
Problem 1.    Let (x1 , x2 , . . .) be a sequence of positive real numbers satisfying                                     = 1. Prove
                                                                                                               n=1
                                                                                                                   2n − 1
that
                                                        ∞ X
                                                          k
                                                        X   xn
                                                                             ≤ 2.
                                                        k=1 n=1
                                                                       k2
                                                          (Proposed by Gerhard J. Woeginger, The Netherlands)

Solution. By interchanging the sums,

                                    ∞ X                                      ∞              ∞
                                      k
                                                                                                     !
                                    X   xn                X xn               X              X 1
                                                   =                     =             xn                .
                                    k=1 n=1
                                            k2         1≤n≤k
                                                             k2              n=1            k=n
                                                                                                k2

Then we use the upper bound
                          ∞               ∞       ∞              
                          X 1             X1     X     1      1         1
                                  ≤          1 =         1 −    1   =
                          k=n
                              k 2
                                    k=n
                                         2
                                        k −4     k=n
                                                      k−2    k+2      n − 12

and get
               ∞ X                  ∞          ∞
                 k
                                                          !        ∞                                ∞
               X   xn               X          X 1                 X                1                 X     xn
                                =         xn                  <              xn ·                =2              = 2.
               k=1 n=1
                          k2        n=1        k=n
                                                     k2            n=1
                                                                                  n − 12              n=1
                                                                                                          2n − 1


Problem 2.        Today, Ivan the Confessor prefers continuous functions f : [0, 1] → R satisfying
                                                                         R1
f (x) + f (y) ≥ |x − y| for all pairs x, y ∈ [0, 1]. Find the minimum of 0 f over all preferred functions.
                                              (Proposed by Fedor Petrov, St. Petersburg State University)
                                 R1
Solution. The minimum of            f is 41 .
                                  0
    Applying the condition with 0 ≤ x ≤ 21 , y = x + 21 we get

                                                   f (x) + f (x + 21 ) ≥ 12 .

By integrating,
                      Z 1                   Z 1/2                                             Z 1/2
                                                       f (x) + f (x + 21 )                            1
                                                                                                        dx = 14 .
                                                                                   
                               f (x) dx =                                              dx ≥           2
                          0                    0                                                 0

   On the other hand, the function f (x) =                    x − 21     satises the conditions because

                                     x − 21 +        1
                                                                   ≤ x − 12 + 21 − y = f (x) + f (y),
                                                             
                  |x − y| =                          2
                                                       −y

and establishes
                    Z 1                   Z 1/2                          Z 1
                                                   1                                        1 1 1
                                                                                 x − 12 dx = + = .
                                                                                      
                          f (x) dx =               2
                                                     −x           dx +
                     0                     0                                 1/2            8 8 4
Problem 3.        Let n be a positive integer, and denote by Zn the ring of integers modulo n. Suppose
that there exists a function f : Zn → Zn satisfying the following three properties:
   (i) f (x) 6= x,
   (ii) f (f (x)) = x,
   (iii) f (f (f (x + 1) + 1) + 1) = x for all x ∈ Zn .
Prove that n ≡ 2 (mod 4).
                       (Proposed by Ander Lamaison Vidarte, Berlin Mathematical School, Germany)

Solution. From property (ii) we can see that    f is surjective, so f is a permutation of the elements in
Zn , and its order is at
                       most 2. Therefore, the permutation f is the product of disjoint transpositions
of the form x, f (x) . Property (i) yields that this permutation has no xed point, so n is be even,
and the number of transpositions is precisely n/2.
    Consider the permutation g(x) = f (x + 1). If g was odd then g ◦ g ◦ g also would be odd. But
property (iii) constraints that g ◦ g ◦ g is the identity which is even. So g cannot be odd; g must
be even. The cyclic permutation h(x) = x − 1 has order n, an even number, so h is odd. Then
f (x) = g ◦ h is odd. Since f is the product of n/2 transpositions, this shows that n/2 must be odd,
so n ≡ 2 (mod 4).


Remark. There exists a function with properties (iiii) for every n ≡ 2 (mod 4).       For n = 2 take f (1) = 2,
f (2) = 1. Here we outline a possible construction for n ≥ 6.
    Let n = 4k + 2, take a regular polygon with k + 2 sides, and divide it into k triangles with k − 1 diagonals.
Draw a circle that intersects each side and each diagonal twice; altogether we have 4k +2 intersections. Label
the intersection points clockwise around the circle. On every side and diagonal we have two intersections;
let f send them to each other.
    This function f obviously satises properties (i) and (ii). For every x we either have f (x + 1) = x or
the eect of adding 1 and taking f three times is going around the three sides of a triangle, so this function
satises property (iii).




Problem 4.     Let k be a positive integer. For each nonnegative integer n, let f (n) be the number
of solutions (x1 , . . . , xk ) ∈ Zk of the inequality |x1 | + ... + |xk | ≤ n. Prove that for every n ≥ 1, we
have f (n − 1)f (n + 1) ≤ f (n)2 .
             (Proposed by Esteban Arreaga, Renan Finder and José Madrid, IMPA, Rio de Janeiro)

Solution 1.    We prove by induction on k . If k = 1 then we have f (n) = 2n + 1 and the statement
immediately follows from the AM-GM inequality.
    Assume that k ≥ 2 and the statement is true for k − 1. Let g(m) be the number of integer
solutions of |x1 | + ... + |xk−1 | ≤ m; by the induction hypothesis g(m − 1)g(m + 1) ≤ g(m)2 holds;
this can be transformed to
                                        g(0)    g(1)   g(2)
                                             ≤       ≤      ≤ ....
                                        g(1)    g(2)   g(3)
   For any integer constant c, the inequality |x1 |+...+|xk−1 |+|c| ≤ n has g n−|c| integer solutions.
                                                                                   

Therefore, we have the recurrence relation
                                           n
                                           X               
                                f (n) =           g n − |c| = g(n) + 2g(n − 1) + ... + 2g(0).
                                           c=−n

It follows that
                                  f (n − 1)     g(n − 1) + 2g(n − 2) + ... + 2g(0)
                                            =                                        ≤
                                     f (n)    g(n) + 2g(n − 1) + ... + 2g(1) + 2g(0)
                               g(n) + g(n − 1) + (g(n − 1) + ... + 2g(0) + 2 · 0)      f (n)
                       ≤                                                          =
                                g(n + 1) + g(n) + (g(n) + ... + 2g(1) + 2g(0))      f (n + 1)
as required.
Solution 2.        We rst compute the generating function for f (n):
            ∞                                       ∞
                                                                                                        !k
            X                          X            X                                     X                   1     (1 + q)k
                  f (n)q n =                                q |x1 |+|x2 |+···+|xk |+c =         q |x|            =            .
            n=0
                                                                                                             1−q   (1 − q)k+1
                                 (x1 ,x2 ,...,xk )∈Zk c=0                                 x∈Z


For each a = 0, 1, . . . denote by ga (n) (n = 0, 1, 2, . . .) the coecients in the following expansion:
                                                                    ∞
                                                       (1 + q)a    X
                                                                 =     ga (n)q n .
                                                      (1 − q)k+1   n=0

So it is clear that ga+1 (n) = ga (n) + ga (n − 1) (n ≥ 1), ga (0) = 1. Call a sequence of positive numbers
g(0), g(1), g(2), . . . good if g(n−1)
                                 g(n)
                                        (n = 1, 2, . . .) is an increasing sequence. It is straightforward to
check that g0 is good:                              
                                              k+n           g0 (n − 1)     n
                                   g0 (n) =             ,              =      .
                                               k               g0 (n)    k+n
If g is a good sequence then a new sequence g 0 dened by g 0 (0) = g(0), g 0 (n) = g(n) + g(n − 1)
(n ≥ 1) is also good:
                          g 0 (n − 1)   g(n − 1) + g(n − 2)   1 + g(n−2)
                                                                  g(n−1)
                               0
                                      =                     =      g(n)
                                                                         ,
                              g (n)       g(n) + g(n − 1)     1 + g(n−1)
where dene g(−1) = 0. Thus we see that each of the sequences g1 , g2 , . . . , gk = f are good. So the
desired inequality holds.

Problem 5.          Let A be a n × n complex matrix whose eigenvalues have absolute value at most 1.
Prove that
                                                                         n
                                                            kAn k ≤          kAkn−1 .
                                                                        ln 2
                                                                                            rn
(Here kBk = sup kBxk for every n × n matrix B and kxk =                                                  |xi |2 for every complex vector
                                                                                             P
                   kxk≤1                                                                          i=1
x ∈ Cn .)
                                 (Proposed by Ian Morris and Fedor Petrov, St. Petersburg State University)

Solution 1.    Let r = kAk. We have to prove kAn k ≤ lnn2 rn−1 .
   As is well-known, the matrix norm satises kXY k ≤ kXk · kY k for any matrices X, Y , and as a
simple consequence, kAk k ≤ kAkk = rk for every positive integer k .
   Let χ(t) = (t − λ1 )(t − λ2 ) . . . (t − λn ) = tn + c1 tn−1 + · · · + cn be the characteristic polynomial
of A. From Vieta's formulas we get
                                                                               
                     X                              X                          n
         |ck | =               λi1 · · · λik ≤                λi1 · · · λik ≤      (k = 1, 2, . . . , n)
                 1≤i <...<i ≤n                  1≤i <...<i ≤n
                                                                               k
                           1       k                              1       k
By the CayleyHamilton theorem we have χ(A) = 0, so
                                            n          n  
              n        n−1
                                            X  n    k
                                                         X  n k
           kA k = kc1 A    + · · · + cn k ≤       kA k ≤      r = (1 + r)n − rn .
                                               k k=1
                                                            k             k=1

Combining this with the trivial estimate kAn k ≤ rn , we have

                                   kAn k ≤ min rn , (1 + r)n − rn ) .
                                                                   

Let r0 = √
         n
           1
           2−1
               ; it is easy to check that the two bounds are equal if r = r0 , moreover

                                                        1              n
                                          r0 =                    <        .
                                                    eln 2/n − 1       ln 2
For r ≤ r0 apply the trivial bound:
                                                                        n n−1
                                   kAn k ≤ rn ≤ r0 · rn−1 <                 r .
                                                                       ln 2
For r > r0 we have
                               n                n       n     n−1       (1 + r)n − rn
                            kA k ≤ (1 + r) − r = r                    ·               .
                                                                            rn−1
                                      n−r  n
Notice that the function f (r) = (1+r)
                                    rn−1
                                          is decreasing because the numerator has degree n − 1 and
all coecients are positive, so
                 (1 + r)n − rn   (1 + r0 )n − r0n                                n
                       n−1
                               <       n−1        = r0 (1 + 1/r0 )n − 1) = r0 <      ,
                     r                r0                                        ln 2

so kAn k < lnn2 rn−1 .
Solution 2. We will use the following facts which are easy to prove:


   • For any square matrix A there exists a unitary matrix U such that U AU −1 is upper-triangular.

   • For any matrices A, B we have kAk ≤ k(A|B)k and kBk ≤ k(A|B)k where (A|B) is the matrix
     whose columns are the columns of A and the columns of B .

   • For any matrices A, B we have kAk ≤ k B   A
                                                   k and kBk ≤ k B
                                                                 A
                                                                     k where B
                                                                             A
                                                                                 is the matrix
                                                                             

     whose rows are the rows of A and the rows of B .

   • Adding a zero row or a zero column to a matrix does not change its norm.
We will prove a stronger inequality
                                               kAn k ≤ nkAkn−1
for any n × n matrix A whose eigenvalues have absolute value at most 1. We proceed by induction
on n. The case n = 1 is trivial. Without loss of generality we can assume that the matrix A is
upper-triangular. So we have                                
                                         a11 a12 · · · a1n
                                       0 a22 · · · a2n 
                                  A= · · · · · ·
                                                             .
                                                        · · ·
                                          0    0 · · · ann
Note that the eigenvalues of A are precisely the diagonal entries. We split A as the sum of 3 matrices,
A = X + Y + Z as follows:
                                                                                            
            a11 0 · · · 0                    0 a12 · · · a1n                 0     0 ··· 0
          0      0 ··· 0      , Y =  0
                                                  0 ··· 0       , Z =  0 a22 · · · a2n  .
                                                                                                
    X=  · · · · · ·      · · ·         · · · · · ·      · · ·        · · · · · ·      · · ·
             0    0 ··· 0                    0     0 ··· 0                   0     0 · · · ann
Denote by A0 the matrix obtained from A by removing the rst row and the rst column:
                                                         
                                           a22 · · · a2n
                                    A0 = · · ·      · · · .
                                            0 · · · ann

We have kXk ≤ 1 because |a11 | ≤ 1. We also have

                                  kA0 k = kZk ≤ kY + Zk ≤ kAk.

Now we decompose An as follows:

                                  An = XAn−1 + (Y + Z)An−1 .

We substitute A = X + Y + Z in the second term and expand the parentheses. Because of the
following identities:
                         Y 2 = 0, Y X = 0, ZY = 0, ZX = 0
only the terms Y Z n−1 and Z n survive. So we have

                                  An = XAn−1 + (Y + Z)Z n−1 .

By the induction hypothesis we have kA0n−1 k ≤ (n − 1)kA0 kn−2 , hence kZ n−1 k ≤ (n − 1)kZkn−2 ≤
(n − 1)kAkn−2 . Therefore

       kAn k ≤ kXAn−1 k + k(Y + Z)Z n−1 k ≤ kAkn−1 + (n − 1)kY + ZkkAkn−2 ≤ nkAkn−1 .
